% s03-3-2-construction.tex —— 3.3.2 与 3.3.3：B 的构造与逐点等价（印刷页 97–102）

\subsection{B 的构造}
\label{sec:b-construction}

本小节利用对的等价 AB 在 $\FactModCat_{x_0}$ 中构造对的态射
\[
\mathrm{B}\colon
(\Sph_G,\Aff_G)\longrightarrow
(\Sph^{\spec}_{\check G},\Aff^{\spec}_{\check G}).
\]

有幺半作用
\[
(\Sph_G,\Aff_G)\otimes (\Whit^{!}(G),\Whit^{!}(\Fl_G))\longrightarrow
(\Whit^{!}(G),\Whit^{!}(\Fl_G)),
\]
\srcpage{97}
它可以写成幺半函子
\[
(\Sph_G,\Aff_G)\longrightarrow
(\Whit^{!}(G),\Whit^{!}(\Fl_G))\otimes
(\Whit^{*}(G),\Whit^{*}(\Fl_G)),
\]
它与
\[
(\Rep(\check G),\QCoh(\check{\mathfrak{n}}/\check B))^{\otimes 2}
\]
等同。
$\Sph_G$ 中因子化（同时也是幺半的）单位在
\[
\Sph_G\to \Rep(\check G)^{\otimes 2}
\]
下的像是因子化代数 $R_{\check G}$，即把 $\check G$ 的正则表示看作 $\check G\times\check G$
的表示。几何上，这就是 $\mathcal{O}_{\LS^{\reg}_{\check G}}$ 沿
\[
\LS^{\reg}_{\check G}\longrightarrow
\LS^{\reg}_{\check G}\times \LS^{\reg}_{\check G}
\]
的推前。

于是得到对的幺半态射
\[
(\Sph_G,\Aff_G)\longrightarrow
R_{\check G}\text{-}\FactMod
\bigl[(\Rep(\check G),\QCoh(\check{\mathfrak{n}}/\check B))^{\otimes 2}\bigr].
\]
由定理~\ref{thm:3.1.5.2}（取 $Y^0=\mathrm{pt}/\check G^2\times X_{\dR}$、
$Y=\mathrm{pt}/\check G\times X_{\dR}$、$S^0=(\check{\mathfrak{n}}/\check B^{\Ran}_{x_0})^2$），
存在态射
\[
(\Sph^{\spec}_{\check G},\Aff^{\spec}_{\check G})\longrightarrow
R_{\check G}\text{-}\FactMod
\bigl[(\Rep(\check G),\QCoh(\check{\mathfrak{n}}/\check B))^{\otimes 2}\bigr],
\]
它在每个分量的最终余连通部分上是等价。

与构造 AB 时一样，我们想说 $\Aff_G$ 紧生成，且其紧生成元被送到最终余连通对象。
然而 $\Aff_G$ 不大可能紧生成。\srcpage{98}可熔性再次救了场：因为
\[
\Aff_{G,\Sph^{\spec}_{\check G}\text{-}\QLisse}\subset \Aff_G
\]
紧生成，且其紧生成元被送到
$R_{\check G}\text{-}\FactMod(\QCoh(\check{\mathfrak{n}}/\check B)^{\otimes 2})^{>-\infty}$。

利用 $\FactRes_{\Rep(\check G)\to \Sph^{\spec}_{\check G}}(\Aff^G)$ 的可熔性，可以把由此得到的
函子
\[
\Aff_{G,\Sph^{\spec}_{\check G}\text{-}\QLisse}\longrightarrow \Aff^{\spec}_{\check G}
\]
提升到整个 $\Aff_G$，从而得到所要的 $\FactModCat_{x_0}$ 中的态射
\[
\mathrm{B}\colon
(\Sph_G,\Aff_G)\longrightarrow
(\Sph^{\spec}_{\check G},\Aff^{\spec}_{\check G}).
\]

\subsection{B 是逐点等价}
\label{sec:b-pointwise}

作相应的因子化限制后，$\Aff_G$ 与 $\Aff^{\spec}_{\check G}$ 作为因子化
$\Rep(\check G)$-模范畴都是可熔的。因此要证 B 是等价，只需证
\[
\mathrm{B}_{x_0}\colon
\Aff_{G,x_0}\longrightarrow \Aff^{\spec}_{\check G,x_0}
\]
是范畴的等价。

在 \cite{DT25} 中，Dhillon 与 Taylor 在单值（monodromic）的 Betti 设定下构造了
（逐点）Bezrukavnikov 等价的一个版本。换言之，它与我们的设定有两个关键差别。
\srcpage{99}
第一，几何侧的叠是
\[
\widetilde{\Fl}_G/I^0=I^0\backslash LG/I^0,
\]
而不是
\[
\Fl_G/I=I\backslash LG/I.
\]
第二，几何侧的层论是带有幂零奇异支撑的 ind-可构造层
\[
\Shv^{\Nilp},
\]
而不是
\[
D\text{-}\mathrm{mod};
\]
相应地，谱侧出现的空间是 Betti 版本。

让我们概述 \cite{DT25} 的构造（更确切地说，是他们的构造在我们的单幂零 de Rham
设定下的类比）。幺半范畴
\[
D\text{-}\mathrm{mod}(\Fl_G)^{(I,\mathrm{ren})}
\]
显然作用在范畴
\[
\Whit^{!}(\Fl_G)
\]
上，后者经 $\mathrm{AB}_{x_0}$ 可等同于
\[
\QCoh(\widetilde{\check N}/\check G).
\]
\srcpage{100}
此外，$D\text{-}\mathrm{mod}(\Fl_G)^{(I,\mathrm{ren})}$ 在 $\Whit^{!}(\Fl_G)$ 上的作用与
\[
\mathrm{bi}\text{-}\Whit(LG)=\Whit^{*}(\Whit^{!}(LG))
\]
在 $\Whit^{!}(\Fl_G)$ 上的作用之间有一个平凡的相容性数据。而后者可以等同于
\[
\QCoh(\check{\mathfrak{g}}/\check G)
\]
经拉回在 $\QCoh(\widetilde{\check N}/\check G)$ 上的作用。于是得到幺半函子
\[
\iota^{!}\colon
D\text{-}\mathrm{mod}(\Fl_G)^{(I,\mathrm{ren})}\longrightarrow
\End_{\QCoh(\check{\mathfrak{g}}/\check G)}
\bigl(\QCoh(\widetilde{\check N}/\check G)\bigr)
\simeq
\QCoh\bigl(\widetilde{\check N}/\check G
\times_{\check{\mathfrak{g}}/\check G}
\widetilde{\check N}/\check G\bigr).
\]
由于 $\QCoh(\widetilde{\check N}/\check G)$ 作为
$D\text{-}\mathrm{mod}(\Fl_G)^{(I,\mathrm{ren})}$-$\QCoh(\check{\mathfrak{g}}/\check G)$
双模的 Morita 性质，$\iota^{!}$ 有全忠实的左伴随 $\iota_!$。此外，$\iota^{!}$ 的核恰由
$D\text{-}\mathrm{mod}(\Fl_G)^{(I,\mathrm{ren})}$ 的无穷连通对象组成，于是得到等价
\[
D\text{-}\mathrm{mod}(\Fl_G)^{(I,\mathrm{ren})}\big/
D\text{-}\mathrm{mod}(\Fl_G)^{(I,\mathrm{ren}),\leq -\infty}
\xrightarrow{\ \sim\ }
\QCoh\bigl(\widetilde{\check N}/\check G
\times_{\check{\mathfrak{g}}/\check G}
\widetilde{\check N}/\check G\bigr).
\]
最后作如下重整化。在谱侧，伪紧对象的全子范畴是
\[
\mathrm{Psc}\bigl(\QCoh(\widetilde{\check N}/\check G
\times_{\check{\mathfrak{g}}/\check G}
\widetilde{\check N}/\check G)\bigr)
=\Coh\bigl(\widetilde{\check N}/\check G
\times_{\check{\mathfrak{g}}/\check G}
\widetilde{\check N}/\check G\bigr).
\]
另一方面，在几何侧有等价
\[
\mathrm{Psc}\bigl(D\text{-}\mathrm{mod}(\Fl_G)^{(I,\mathrm{ren})}\bigr)
\simeq
\mathrm{Psc}\bigl(D\text{-}\mathrm{mod}(\Fl_G)^{(I,\mathrm{ren})}\big/
D\text{-}\mathrm{mod}(\Fl_G)^{(I,\mathrm{ren}),\leq -\infty}\bigr),
\]
且 $D\text{-}\mathrm{mod}(\Fl_G)^{(I,\mathrm{ren})}$ 中的伪紧对象恰是紧对象（即
$D\text{-}\mathrm{mod}(\Fl_G)^{I}$ 中在两侧遗忘到 $D\text{-}\mathrm{mod}(\Fl_G)$ 后都紧的对象）。
于是得到等价
\srcpage{101}
\[
\mathrm{B}^{\mathrm{classical}}\colon
D\text{-}\mathrm{mod}(\Fl_G)^{(I,\mathrm{ren})}
\xrightarrow{\ \sim\ }
\IndCoh\bigl(\widetilde{\check N}/\check G
\times_{\check{\mathfrak{g}}/\check G}
\widetilde{\check N}/\check G\bigr).
\]

现在的任务是：把因子化函子 B 在 $x_0$ 处的纤维 $\mathrm{B}_{x_0}$ 与
$\mathrm{B}^{\mathrm{classical}}$ 等同起来。为此只需说明
\[
\Aff_{G,x_0}
\xrightarrow{\ \mathrm{B}_{x_0}\ }
\Aff^{\spec}_{\check G,x_0}
\xrightarrow{\ \Psi_{x_0}\ }
\QCoh_{\mathrm{co}}\Bigl(\check{\mathfrak{n}}/\check B^{\Ran}_{x_0}
\times_{(\LS^{\mer}_{\check G})_{\Ran_{x_0}}}
\check{\mathfrak{n}}/\check B^{\Ran}_{x_0}\Bigr)^{x_0}
\]
与
\[
D\text{-}\mathrm{mod}(\Fl_G)^{(I,\mathrm{ren})}
\xrightarrow{\ \mathrm{B}^{\mathrm{classical}}\ }
\IndCoh\bigl(\widetilde{\check N}/\check G
\times_{\check{\mathfrak{g}}/\check G}
\widetilde{\check N}/\check G\bigr)
\xrightarrow{\ \Psi\ }
\QCoh\bigl(\widetilde{\check N}/\check G
\times_{\check{\mathfrak{g}}/\check G}
\widetilde{\check N}/\check G\bigr)
\]
中的 $\iota_!$ 相合。由构造，我们已经有
\[
\Aff_{G,x_0}
\xrightarrow{\ \mathrm{B}_{x_0}\ }
\Aff^{\spec}_{\check G,x_0}
\xrightarrow{\ \Psi_{x_0}\ }
\QCoh_{\mathrm{co}}\Bigl(\check{\mathfrak{n}}/\check B^{\Ran}_{x_0}
\times_{(\LS^{\mer}_{\check G})_{\Ran_{x_0}}}
\check{\mathfrak{n}}/\check B^{\Ran}_{x_0}\Bigr)^{x_0}
\]
\begin{equation}
\longrightarrow
\QCoh\bigl(\check{\mathfrak{n}}/\check B^{\Ran}_{x_0}
\times \check{\mathfrak{n}}/\check B^{\Ran}_{x_0}\bigr)^{x_0}
\label{eq:3.3.6}
\end{equation}
与
\begin{equation}
D\text{-}\mathrm{mod}(\Fl_G)^{(I,\mathrm{ren})}
\xrightarrow{\ \mathrm{B}^{\mathrm{classical}}\ }
\IndCoh\bigl(\widetilde{\check N}/\check G
\times_{\check{\mathfrak{g}}/\check G}
\widetilde{\check N}/\check G\bigr)
\xrightarrow{\ \Psi\ }
\QCoh\bigl(\widetilde{\check N}/\check G
\times_{\check{\mathfrak{g}}/\check G}
\widetilde{\check N}/\check G\bigr)
\label{eq:3.3.7}
\end{equation}
\begin{equation}
\longrightarrow
\QCoh\bigl(\widetilde{\check N}/\check G\times
\widetilde{\check N}/\check G\bigr)
\label{eq:3.3.8}
\end{equation}
的等同（也就是说，由两个作用定义的到
$\End(\QCoh(\widetilde{\check N}/\check G))$ 的函子相同，因为后一个作用是用
$\mathrm{AB}_{x_0}$ 构造的）。
\srcpage{102}

因此，剩下的全部工作是说明：由 $\Psi_{x_0}\circ \mathrm{B}_{x_0}$ 给出的
$\Aff_{G,x_0}$ 在 $\QCoh(\widetilde{\check N}/\check G)$ 上的作用与
$\QCoh(\check{\mathfrak{g}}/\check G)$ 的作用之间的相容性数据，与平凡的相容性数据同构。
但这可由以下事实推出：对
\[
(\Rep(\check G),\QCoh(\check{\mathfrak{n}}/\check B))
\]
带有对
\[
\bigl(\QCoh(\LS^{\mer}_{\check G}),\QCoh(\LS^{\mer}_{\check G})\bigr)
\]
的幺半作用，而它与 $(\Sph_G,\Aff_G)$ 的幺半作用之间有平凡的相容性数据；并且限制到
$x_0$ 时，前一个作用等同于 $\QCoh(\LS_{\check G}(D^\times_{x_0}))$ 经拉回
\[
\widetilde{\check N}/\check G\longrightarrow \LS_{\check G}(D^\times_{x_0})
\]
（该映射经 $\check{\mathfrak{g}}/\check G$ 分解）在 $\QCoh(\widetilde{\check N}/\check G)$
上的作用。这就完成了定理~\ref{thm:B} 的证明。
